\subsection{Recovery Manager: Preassigned Takeover}
\label{sec:parallel-recovery}
\label{sec:recovery-manager}

Recovery Manager (2) decides where each request will resume before
failure occurs. When request $r$ arrives and the Router selects source
DP rank $g_r$, the Manager records a target assignment $p_r\ne g_r$
in control state outside the source's failure domain. The same $p_r$
is both the retry destination and the DP rank that will hold any
Peer-HBM backup of $r$. A simple mapping for $n\ge2$ DPs is
$p_r=(g_r+1)\bmod n$; per-request assignments may distribute load across
several targets. The assignment does not reserve HBM: the KV Master
chooses whether to request Peer protection, and the target Adapter
controls its capacity. A TP group restores compatible shards as one DP.

\paragraph{Activate the recorded route after failure.}
The Manager combines membership, device errors, and GPU progress
reported by Adapters; a process heartbeat alone does not prove GPU
progress. After confirming and isolating a failed source, it obtains
request token prefixes, commit frontiers, model/layout IDs, and replica
locations from surviving control state, then submits retries to their
recorded targets. Available Peer replicas are therefore local to takeover.
DRAM or recomputation supplies the remaining state. Control-state
survival, fencing, failure-domain separation, and retry deduplication
are runtime contracts. A failed target requires a new valid assignment;
copies on the old target cannot be treated as local to the new target.

\paragraph{Recover a dependency-closed frontier.}
A surviving authority assigns a monotonically increasing source epoch;
copies, directory commits, and output acknowledgements reject stale epochs.
Before acknowledging a generated token, persist its request ID, position,
token ID, and epoch in a surviving ordered log. Prompt and accepted tool
inputs are likewise recoverable. KV may lag this log; recovery reconstructs
the acknowledged suffix and suppresses duplicate delivery. External tool
side effects require the orchestrator's own idempotency contract.

Publish a checkpoint manifest only after its full dependency closure is
committed: compatible model, weights and layout, all required layers and TP
shards, sealed token intervals, selector state, and the matching token-log
prefix. A frontier is the latest valid common manifest, not the maximum
of independently updated object locations. Peer and DRAM may commit in
either order; each manifest entry must resolve to a matching surviving
replica. An incomplete new manifest leaves the older one authoritative;
without one, replay from recoverable inputs. Recovery validates versions
and closure again after fencing. These checks define the minimum contract;
their concurrent implementation remains to be validated.

\paragraph{Share resources without serializing retries.}
Retries enter independently and advance on their own layer-readiness
events. Slack determines protection, not a request-level completion
barrier. The engine reserves a positive dispatch share for feasible recovery
work while continuing foreground batching. Within that share, ready admitted
retries receive round-robin service in bounded prefill chunks or decode steps;
eligible waiting retries use oldest-first admission subject to physical space.
Blocked layers yield rather than holding a request-level barrier. These are
required scheduler interfaces, not a claim of support in an unmodified engine.
Each retry moves through waiting-for-space, waiting-for-data, ready, and
running states. Reserve its bounded activation/KV workspace and transfer
buffers before admitting an execution chunk. The batch coordinator admits
compatible ready work and gives every TP rank the same batch, layer, and
collective order. A newly discovered miss parks that work only at a legal
kernel boundary, retaining and charging its activation buffer; it can split
from a batch only if the backend supports that split. Otherwise the whole
batch waits and the profile charges the resulting interference. Readiness
does not authorize arbitrary layer reordering or unmatched collectives.
Memory reservations precede I/O leases; release queue/directory locks before
waiting on a copy or collective, and reserve demand-transfer progress space.
Preassignment reserves neither compute nor retry workspace. Profiles
include target foreground batches, available KV/workspace, readmission
queues, shared links and failure domains. If the target has no feasible
working set, queue the retry or remap through the Manager with a new
placement check; never assume that local backup alone permits execution.

Transfers contend at every shared link and at target HBM. Normal-state
backup traffic has a configured aggregate cap (default 20\% of link
capacity). During recovery or contraction, configure positive service
shares for demand fetch and required writeback while retaining a foreground
share; active shares sum to at most one per bottleneck and unused service
can be borrowed. Actual demand precedes speculative preloads within recovery.
The 20\% traffic cap is not a 20\% bound on foreground latency or overhead.
Byte-based deficit round robin across eligible requests and a bounded
outstanding-transfer count prevent one long context from monopolizing
the recovery share. Migration writeback receives its own progress share,
so shrinking does not depend on all decode requests ending. Service switches
at supported DMA/kernel boundaries, not at an assumed arbitrary instant.

The same sharing policy and foreground/retry load must enter the cost
profiles. With positive available service, finite contenders, and bounded
transfer quanta, eligible work progresses; this is not a latency guarantee
under overload. No retry must complete before another can begin merely
because its placement priority was higher.
